\chapter{Personal Credit Strategy}

I currently have no debt and no credit card---a clean position, but also no credit history. Vietnam's CIC (Credit Information Center) system records whether borrowers take out credit products and repay them consistently. Banks and lenders check this record when evaluating any loan application. A borrower with no record at all is treated with caution, because there is no evidence of how they handle debt. Building a score strong enough to support a 10 billion VND mortgage application in 2040 requires a minimum of three to five years of consistent credit card use---which means the first card needs to be opened well before the mortgage is needed.

\section{Building a Credit History}

Three steps, in order. First, get my first credit card in 2029---either a Techcombank Signature or ABBank Visa, with a limit roughly equal to one month's gross salary (50,000,000 to 70,000,000~VND at BCG starting pay). This opens a credit product that generates a repayment record on the CIC system. Second, keep spending below 30\% of the available limit at all times. Banks read low utilization as a sign that I borrow by choice, not out of necessity---a borrower consistently at 80\% of their limit suggests financial stress, while one consistently at 20 to 30\% signals control. Third, set up automatic full-balance payment on the monthly due date. This removes any risk of late-payment marks on the CIC record and requires zero ongoing management after setup.

Two behaviors that significantly damage a CIC score: missing a payment deadline (even by one day) and carrying a revolving balance that accumulates interest. Both are completely avoidable through auto-debit. A clean, unbroken 11-year payment record from 2029 to 2040 is the single strongest signal a Vietnamese bank can receive from a retail mortgage applicant.

\textbf{Practical Example.} The Vietcombank VCB Visa Classic is a strong first card: initial limit of 10,000,000 to 50,000,000~VND for verified-salary applicants, 45-day interest-free period, auto-debit payment available through VCB digital banking, and accepted at virtually every merchant and online platform in Vietnam. Spending under 15,000,000~VND/month on a 50,000,000~VND limit and paying automatically builds a clean CIC record within 12 months of activation. Starting in 2029 gives 11 years of documented credit history before the 2040 mortgage application.

\section{Long-Term Debt Strategy}

For the Thanh Xu\^{a}n house in 2040: total price 20,000,000,000~VND, split 50-50 between personal funds and bank debt. The 10,000,000,000~VND bank loan runs over 240 months at an estimated 8.5\% per annum, producing a monthly payment of approximately \textbf{86,780,000~VND}. At a projected income of 300,000,000~VND/month as BCG Partner, this is a Debt Service Ratio of 28.9\%---below the 35\% safe lending ceiling used by Vietnamese banks.

The 50\% Loan-to-Value ratio provides a meaningful equity buffer. If property prices fall during the 20-year repayment period, the outstanding loan balance remains below the property value as long as prices do not drop more than 50\% from the 2040 purchase level---an extremely unlikely scenario for a well-located house in an urban district of Hanoi. The combination of a strong CIC record, a verified high income at BCG salary levels, and a 50\% down payment means the mortgage application in 2040 should be approved at competitive rates. The credit-building work that starts in 2029 is what makes the 2040 application straightforward rather than difficult.
